\section*{Principio de Equivalencia}
\underline{Principio de Equivalencia Débil}:
\begin{align*}
&\text{Para partículas con masa es imposible discernir si una partícula está sometida a los}\\
&\text{efectos de gravedad uniforme o aceleración uniforme}\Rightarrow \vec{g}\!=\!\vec{a} \; \; m_{\scalebox{0.8}{\text{gravitatoria}}}\!=\!  m_{\scalebox{0.8}{\text{inercial}}}
\end{align*}
\underline{Principio de Equivalencia Fuerte (o de Einstein)}:
\begin{align*}
&\text{Para toda partícula y toda fuerza es imposible discernir si una partícula está}\\
&\text{sometida a los efectos de  de gravedad uniforme o aceleración uniforme.}
\end{align*}
\underline{Teorema}:
\begin{align*}
    &\text{En la vecindad de cualqier punto $P$ del espacio-tiempo, siempre podemos escoger}\\
    &\text{coordenadas de manera que: } g_{\mu \nu }= 
    \eta_{\mu\nu} \text{ en el punto $P$.}\;\;\; \eta_{\mu\nu}= \text{diag}(-1,1,1,1) 
    \end{align*}
    \vspace{1pt}
\begin{align*}
\blacktriangleright
\text{ S.R. Inercial (Minkowski)} \equiv \text{Sistema de Referencia en caída libre.}
\end{align*}